\chapter{超電荷、閉包、超空間}

\lectureoverview{超対称性が難しいのは、その規則が多いからではなく、その規則が通常の物理学では別々に保たれている考えを結びつけるからである。まずグラスマン変数から始める。その奇要素は反可換であり、その関数は有限個の項の後で終わる。次にその微分代数が奇対称性生成子の模型を与える。ボソン振動子とフェルミオン振動子は最も単純な超電荷を与える。それらの反交換子はハミルトニアンであるから、ボソンとフェルミオンを交換することは時間並進と不可分である。同じ代数は、超空間と呼ばれる拡大された時空のグラスマン方向の並進によって幾何学的に表現できる。}

\section{先に進む前に立ち戻らねばならない理由}

超対称性は物理学で最も抽象的な構成の一つである。それは対称性であるから回転や並進に似ているが、座標の意味を拡大することなしには通常の時空の中に収まらない。その自然なパラメータはグラスマン量である。すなわち、加算と乗算ができるが、可換ではなく反可換である対象である。

それゆえ、最終的な代数に直接取り組むことは役に立たないであろう。まず奇数の算術を使えるようにし、次に通常の連続対称性がどのように組織されるかを思い出し、その後に初めて二つの考えを組み合わせる。道のりは長いが、各段階が次の段階を準備する。

\section{実用的な形でのグラスマン代数}

\subsection{奇要素と偶要素}

\(\epsilon_i\)を独立なグラスマン生成子とする。それらは
\begin{equation}
\epsilon_i\epsilon_j=-\epsilon_j\epsilon_i,
\qquad
\{\epsilon_i,\epsilon_j\}=0.
\end{equation}
を満たす。\(i=j\)と置くと
\begin{equation}
2\epsilon_i^2=0,
\qquad
\epsilon_i^2=0,
\end{equation}
が実数または複素数上で得られる。生成子はゼロではない。その二乗だけがゼロである。

各同次要素はパリティをもつ。奇数個のグラスマン生成子を含む積は奇であり、偶数個を含む積は偶である。
\begin{align}
\text{odd}\times\text{odd}&=\text{even},\\
\text{odd}\times\text{even}&=\text{odd},\\
\text{even}\times\text{even}&=\text{even}.
\end{align}
偶要素は同次要素と可換である。二つの奇要素は反可換である。したがって、\(n\)が偶であるとき
\begin{equation}
n\epsilon_i=\epsilon_i n,
\qquad
\epsilon_i\epsilon_j=-\epsilon_j\epsilon_i,
\end{equation}
が成り立つ。

偶要素は通常の実数や複素数である必要はない。例えば積
\(\epsilon_1\epsilon_2\)
は偶であり、かつ冪零である。通常の数は代数のスカラー部分を形成し、偶の積はさらに可換な要素を与える。

\begin{classroomqa}
\audiencequestion{グラスマン代数とグラスマン数の間には区別があるのか、それとも同じものなのか。\lecturetimestamp{00:06:45}}
\lecturerresponse{要素がグラスマン数であり、代数はそれらの要素をその加算と乗算の規則とともに集めたものである。したがって、同じ対象について大まかに語ることはできるが、操作されるものとそれを操作する規則とを区別することは有用である。}
\end{classroomqa}

\subsection{関数は終わる}

一つのグラスマン変数の関数は二つの冪しかもち得ない。
\begin{equation}
f(\epsilon)=a+b\epsilon.
\end{equation}
\(\epsilon^2\)の項はない。二つの独立な生成子があるとき、
\begin{equation}
f(\epsilon_1,\epsilon_2)
=a+b_1\epsilon_1+b_2\epsilon_2+c\epsilon_1\epsilon_2.
\end{equation}
展開は両方の生成子が現れた後で止まる。\(n\)個の独立な生成子に対して、順序づけられた単項式は\(2^n\)次元の基底を形成する。
\begin{equation}
1,\quad
\epsilon_i,\quad
\epsilon_i\epsilon_j,\quad
\ldots,\quad
\epsilon_1\epsilon_2\cdots\epsilon_n.
\end{equation}
\editorialnote{有限な\(2^n\)次元の言明は、一生成子および二生成子の例の背後にある一般的なパターンである。無限のグラスマン代数も定義できるが、ここで用いられるすべての多項式は有限個の生成子しか含まない。}

二つの実生成子が複素対\(\theta,\bar\theta\)に組み合わされるとき、一般の関数は、選ばれた順序で
\begin{equation}
f(\theta,\bar\theta)
=a+\bar b\,\theta+\bar\theta\,b+c\,\bar\theta\theta.
\end{equation}
と書ける。バーはここで用いられる複素対の記法を示す。最後の順序を逆にすると符号が変わる。
\begin{equation}
\bar\theta\theta=-\theta\bar\theta.
\end{equation}

通常は関数全体に確定したパリティを割り当てると便利である。\(f\)が偶ならば、\(a\)と\(c\)は偶であり、\(b\)と\(\bar b\)は奇である。\(f\)が奇ならば、それらの割り当ては逆になる。係数自体は、展開している特定の変数に依存しない限り、より大きなグラスマン代数に属していてもよい。

\begin{classroomqa}
\audiencequestion{グラスマン変数を三角関数やその他の解析関数に挿入できるか。\lecturetimestamp{00:11:06}}
\lecturerresponse{できるが、冪級数は直ちに終わる。一つの奇生成子に対して、
\[
\sin\epsilon=\epsilon,\qquad
\cos\epsilon=1,\qquad
e^\epsilon=1+\epsilon.
\]
おなじみの名前は生き残るが、関数が通常の解析的複雑さを発展させるには冪が少なすぎる。}
\end{classroomqa}

\section{微分は符号の記帳である}

我々は左グラスマン微分を用いる。その定義的な局所規則は
\begin{equation}
\partial_\theta a=0,
\qquad
\partial_\theta(\theta B)=B,
\end{equation}
であり、ここで\(a\)と\(B\)は\(\theta\)に依存しない。微分する前に、関連する\(\theta\)を左へ移動させる。別の奇因子を通るたびに符号が変わる。

例えば、\(b\) が奇数のとき、
\begin{equation}
b\theta=-\theta b,
\qquad
\partial_\theta(b\theta)=-b.
\end{equation}
微分に不思議なことは何も起きていない。マイナス符号は微分の前に、奇数の因子を並べ替えたときに生じたのである。

\subsection{2変数の計算}


\begin{equation}
f(\theta,\bar\theta)
=a+\bar b\,\theta+\bar\theta\,b+c\,\bar\theta\theta.
\end{equation}
に戻る。まず \(f\) が偶数であると仮定する。すると \(\bar b\) は奇数であり、\(c\) は偶数である。\(\theta\) を含む項を並べ替えると、
\begin{align}
\bar b\,\theta&=-\theta\bar b,\\
c\,\bar\theta\theta&=-\theta c\,\bar\theta.
\end{align}
したがって
\begin{equation}
\partial_\theta f=-\bar b-c\bar\theta.
\end{equation}

\(f\) が奇数ならば、\(\bar b\) は偶数であり、\(c\) は奇数である。最初の項は奇数の係数を越えないが、最後の \(\theta\) は2つの奇数因子 \(c\) と \(\bar\theta\) を越える。2つの符号は打ち消し合う：
\begin{equation}
\partial_\theta f=\bar b+c\bar\theta.
\end{equation}

\begin{classroomqa}
\audiencequestion{偶数の例で、なぜ生き残る両方の項がマイナス符号を持つのか。\lecturetimestamp{00:16:49}}
\lecturerresponse{最初の符号は、\(\theta\) を奇数の係数 \(\bar b\) の前を通過させたことから生じる。最後の項では、\(c\) は偶数であるが、\(\theta\) は奇数の変数 \(\bar\theta\) を通過しなければならない。それぞれの並べ替えは1つの奇数因子を越えるので、それぞれが1つのマイナス符号を生じる。}
\end{classroomqa}

答えは順序の規約に依存するが、数学は依存しない。\(c\bar\theta\theta\) ではなく \(c\theta\bar\theta\) から始めていたならば、表示される係数は符号だけ異なり、微分の公式もそれに対応して変わるであろう。

\begin{classroomqa}
\audiencequestion{最後の項を \(\theta\) と \(\bar\theta\) を逆にして書くことはできたか、そしてその結果生じる符号にはより深い意味があるのか。\lecturetimestamp{00:18:50}}
\lecturerresponse{どちらの順序でも書けるが、2つの奇数変数を交換すると符号が変わる。その符号を係数 \(c\) に吸収してもよい。この段階では効果は帳簿付けである：順序は重要であり、パリティは変わらず、符号の変化が奇数な対象を偶数に変えることは決してない。物理的な意味は、同じ帳簿付けがフェルミオン統計と超対称性を強制するときに現れる。}
\end{classroomqa}

\begin{classroomqa}
\audiencequestion{係数 \(a,b,c\) は単に複素数体から取られたものか。\lecturetimestamp{00:20:47}}
\lecturerresponse{それらが必要とされるパリティが偶数のときは通常の複素数であってよいが、選択された \(\theta,\bar\theta\) 以外の生成子を含むこともある。関心のある代数は通常、すべての複素スカラーと、選ばれた奇数の生成子およびその積を含む。}
\end{classroomqa}

\begin{classroomqa}
\audiencequestion{偶関数と奇関数の導関数は反対に見える。これは一般的な対称性だろうか？ \lecturetimestamp{00:22:31}}
\lecturerresponse{それだけでは言えない。符号は各係数のパリティと、交差する奇因子の個数から機械的に従う。二つの答えがたまたま反対になることもあるが、信頼できる規則はもっと単純である。すなわち、二つの奇な対象を交換するごとにマイナス符号が一つ寄与する。}
\end{classroomqa}

\subsection{座標と微分の演算子としての扱い}

示唆に富む通常の類似がある。\(x\) による掛け算と \(x\) による微分を任意の関数 \(g(x)\) に作用させるとする。すると
\begin{align}
\partial_x\!\left(xg\right)&=g+xg',\\
x\partial_xg&=xg',
\end{align}
したがって
\begin{equation}
[\partial_x,x]g=g.
\end{equation}
これがすべての \(g\) について成り立つので、
\begin{equation}
[\partial_x,x]=1.
\end{equation}
\(p=-i\partial_x\) とすると、これは \(\hbar=1\) の単位でのおなじみの位置-運動量代数である。

グラスマン類似では反交換子を用いる。
\begin{equation}
f(\theta)=a+\theta b,
\end{equation}
で \(a\) は偶数、\(b\) は奇数とし、したがって \(f\) は偶数とする。まず、
\begin{equation}
\theta\partial_\theta f=\theta b.
\end{equation}
逆の順序では、
\begin{align}
\partial_\theta(\theta f)
&=\partial_\theta(\theta a+\theta^2b)\\
&=a.
\end{align}
2つの順序を加えると元の関数が得られる：
\begin{equation}
\left(\partial_\theta\theta+\theta\partial_\theta\right)f
=a+\theta b=f.
\end{equation}
したがって
\begin{equation}
\boxed{\{\partial_\theta,\theta\}=1.}
\end{equation}
微分 \(\partial_\theta\) 自体が奇数である。だからこそ、奇数座標との正準関係が反交換子となるのである。

\begin{classroomqa}
\audiencequestion{明確なパリティを割り当てるのではなく、偶数部分と奇数部分の両方を含む関数を用いてもよいか。\lecturetimestamp{00:31:05}}
\lecturerresponse{そのような非斉次な要素を形成することはできるが、計算は通常、斉次な偶数および奇数な対象に整理される。明確なパリティがあれば、次数付き積のすべての符号が曖昧でなくなる。混合した対象は、その偶数部分と奇数部分に分けることで扱われる。}
\end{classroomqa}

\begin{classroomqa}
\audiencequestion{通常の交換子の計算におけるテスト関数 \(g(x)\) は奇数であるべきだったのか。\lecturetimestamp{00:31:45}}
\lecturerresponse{いいえ。それは通常の可換な座標の通常の関数であった。それは、グラスマン対 \(\partial_\theta,\theta\) で議論を繰り返す前に、演算子の恒等式 \([\partial_x,x]=1\) がどのように証明されるかを示すために導入されただけである。}
\end{classroomqa}

\begin{classroomqa}
\audiencequestion{独立なグラスマン生成子の数に限界はあるか。\lecturetimestamp{00:32:16}}
\lecturerresponse{代数的な上限はない。\(n\) 個の生成子に対して基底は \(2^n\) 個の単項式を持つので、帳簿付けは急速に増大するが、構成自体は任意の有限な \(n\) に対して機能する。}
\end{classroomqa}

\subsection{積分}

1つのグラスマン変数にわたる積分は次で固定される：
\begin{equation}
\int d\theta\,1=0,
\qquad
\int d\theta\,\theta=1.
\end{equation}
したがって、\(f(\theta)=a+\theta b\) に対して、
\begin{equation}
\int d\theta\,f(\theta)=b=\partial_\theta f(\theta).
\end{equation}
よって積分と微分は同じ成分抽出を行う。これがベレジン積分である。その規格化は、並進不変性とグラスマン部分積分と両立する唯一の単純な選択である。\editorialnote{成分規則は、グラスマン積分と微分が同じ操作であるという講義の主張を精密化する。}

\section{通常の対称性を基準として}

エルミート生成子 \(G_i\) をもつ連続ユニタリ群を考える。単位元の近くでは、
\begin{equation}
U(\alpha)=\mathbf 1+i\alpha^iG_i+O(\alpha^2),
\end{equation}
ここで \(\alpha^i\) は通常の可換なパラメータである。無限小の乗法則はリー代数
\begin{equation}
[G_i,G_j]=i f_{ij}{}^{k}G_k.
\end{equation}
によって符号化される。定数 \(f_{ij}{}^k\) は構造定数である。閉性とは、任意の2つの生成子を可換にすると、すでに代数に含まれる生成子の線形結合が返ることを意味する。

回転に対しては、
\begin{equation}
[L_i,L_j]=i\varepsilon_{ijk}L_k.
\end{equation}
3つのエルミート成分は非エルミートのはしご演算子
\begin{equation}
L_\pm=L_x\pm iL_y,
\end{equation}
に組み合わせることができ、これらは磁気量子数 \(m\) を1だけ上げ下げする。複素結合は、基礎となる物理的生成子のエルミート性を変えることなく有用である。

一般の基底では、\(f_{ij}{}^k\) の最初の2つの添字だけが反対称でなければならない。不変な正規直交計量をもつコンパクト代数では、最後の添字を下げ、回転の場合と同様に \(f_{ijk}\) が完全反対称となる規約を選ぶことができる。\editorialnote{この限定は、普遍的に反対称な交換子と、なじみ深い回転基底で利用可能なより強い性質とを区別する。}

\(G_i\) が時間に依存しない系の実際の対称性を生成するためには、
\begin{equation}
[H,G_i]=0.
\end{equation}
同値に、\(G_i\) はすべてのエネルギー固有空間をそれ自身へ写す。時間依存が陽にないとき、ハイゼンベルク方程式は
\begin{equation}
\frac{dG_i}{dt}=i[H,G_i]=0,
\end{equation}
を与えるので、生成子は保存電荷である。

通常の観測量は偶である。フェルミオン場は奇であるが、孤立した奇の場それ自体は測定可能な観測量ではない；物理的演算子は偶数個のフェルミオン因子の総数を含む。超対称性は生成子のパリティを変えるのであり、観測結果に要求されるパリティを変えるのではない。

すべての場合に対する明快な記法は次数付き交換子
\begin{equation}
[X,Y]_{\mathrm{gr}}
=XY-(-1)^{|X||Y|}YX.
\end{equation}
である。これは、どちらかの演算子が偶なら通常の交換子であり、両方が奇なら反交換子である。

\section{最小の超対称振動子}

\subsection{1つのボソンモードと1つのフェルミオンモード}

生成消滅演算子 \(a_+\) と \(a_-\) をもつ単一のボソンモードを取る：
\begin{equation}
[a_-,a_+]=1,
\qquad
[a_+,a_+]=[a_-,a_-]=0.
\end{equation}
その数演算子は
\begin{equation}
N_B=a_+a_-,
\end{equation}
であり、その占有数は \(0,1,2,\ldots\) でありうる。

単一のフェルミオンモードに対して、\(c_+\) と \(c_-\) を導入する：
\begin{equation}
\{c_-,c_+\}=1,
\qquad
\{c_+,c_+\}=\{c_-,c_-\}=0.
\end{equation}
したがって
\begin{equation}
c_+^2=c_-^2=0,
\qquad
N_F=c_+c_-,
\end{equation}
であり、占有数は \(0\) または \(1\) のみである。ボソン演算子は偶、フェルミオン演算子は奇であり、すべての \(a_\pm\) はすべての \(c_\pm\) と可換である。

\editorialnote{冪零性はフェルミオンモードを1つの占有状態に制限する。ボソンモードの占有は制限されない；教室での議論における一時的な反対の文言は、振動子代数そのものによって解決される。}

ボソンとフェルミオンが同じ静止エネルギー \(m\) をもつと仮定し、
\begin{equation}
Q=\sqrt m\,a_+c_-,
\qquad
\bar Q=Q^\dagger=\sqrt m\,a_-c_+.
\end{equation}
と定義する。演算子 \(Q\) は1つのフェルミオンを消し、1つのボソンを生成する。その共役は1つのボソンを消し、1つのフェルミオンを生成する。それぞれが1つのフェルミオン演算子を含むので、奇である。

数状態の上で、それらの作用は陽に
\begin{align}
Q\ket{n,1}
&=\sqrt{m(n+1)}\,\ket{n+1,0},
&
Q\ket{n,0}&=0,\\
\bar Q\ket{n,0}
&=\sqrt{mn}\,\ket{n-1,1},
&
\bar Q\ket{n,1}&=0.
\end{align}
である。\editorialnote{数状態の公式は、講義で述べられた粒子置換作用の直接的な規格化された形である。}

\subsection{反交換子はエネルギーで閉じる}

リー代数を見つけることの次数付き類似は、
\begin{align}
\{Q,\bar Q\}
&=Q\bar Q+\bar Q Q\\
&=m\left(a_+c_-a_-c_+
        +a_-c_+a_+c_-\right).
\end{align}
を計算することである。\(a\) は偶なので、符号なしで \(c\) を通り抜ける：
\begin{equation}
\{Q,\bar Q\}
=m\left(a_+a_-c_-c_+
        +a_-a_+c_+c_-\right).
\end{equation}
ここで
\begin{equation}
c_-c_+=1-c_+c_-,
\qquad
a_-a_+=a_+a_-+1.
\end{equation}
を用いる。すると
\begin{align}
\{Q,\bar Q\}
&=m\left[
N_B(1-N_F)+(N_B+1)N_F
\right]\\
&=m(N_B+N_F).
\end{align}
ライブの並べ替えの間に生じた一時的な余分な因子2は消える：一方の寄与が \(N_B\) を与え、他方が \(N_F\) を与える。

等しい質量をもつ静止粒子に対して、ハミルトニアンは
\begin{equation}
H=m(N_B+N_F).
\end{equation}
である。したがって
\begin{equation}
\boxed{\{Q,\bar Q\}=H.}
\end{equation}

\begin{figure}[H]
\centering
\includegraphics[width=0.62\textwidth]{lecture_05_figure_02.png}
\lectureframe{00:55:46.300}
\caption{黒板で分離された中心的な閉関係。中括弧と \(Q,\bar Q,H\) は見えており、コンマは標準的な反交換子記法によって補われる。}
\end{figure}

エネルギーの加法的ゼロは、空状態が \(H=0\) をもつように選ばれている。超対称振動子ではこれが自然な規約である：ボソンとフェルミオンの零点寄与が相殺する。\editorialnote{講義は全静止エネルギーを直接扱い、別々の零点項を表示しない；述べられた規約がその選択を明示する。}

ハミルトニアンは偶である。それはまた対称性の生成子でもある。すなわち、時間並進を生成する。したがって、偶の時間並進生成子を含めない限り、奇の生成子はそれらの間で閉じない。

\begin{classroomqa}
\audiencequestion{この関係のうち偶であるのはどの要素か。ハミルトニアンか。\lecturetimestamp{00:58:24}}
\lecturerresponse{はい。エネルギーは通常の観測量であるから、\(H\)は偶である。超電荷\(Q\)は奇であり、二つの奇の量の積は偶である。したがって、\(\{Q,\bar Q\}\)は\(H\)に等しくなりうる。奇の生成子と偶の生成子は反交換子ではなく交換子によって関係づけられる。}
\end{classroomqa}

\subsection{代数の完成}

同一パリティの関係は直ちに得られる：
\begin{equation}
\{Q,Q\}=2Q^2=0,
\qquad
\{\bar Q,\bar Q\}=2\bar Q^2=0,
\end{equation}
なぜなら\(Q^2\)は\(c_-^2\)を含み、\(\bar Q^2\)は\(c_+^2\)を含むからである。

残るのは\(H\)との関係である。
\begin{equation}
H\ket{E}=E\ket{E}.
\end{equation}
ボソンとフェルミオンの質量が等しいならば、\(Q\)はエネルギーを変えずに一方を他方に置き換える。したがって、\(Q\ket{E}\)はゼロであるか、同じ固有値をもつ別の固有状態である：
\begin{equation}
H Q\ket{E}=E Q\ket{E}.
\end{equation}
一方、
\begin{equation}
QH\ket{E}=Q(E\ket{E})=E Q\ket{E}.
\end{equation}
差をとると、
\begin{equation}
[H,Q]\ket{E}=0.
\end{equation}
エネルギー固有状態が完全な基底をなすならば、これは演算子の恒等式として成り立つ：
\begin{equation}
[H,Q]=0,
\qquad
[H,\bar Q]=0.
\end{equation}

\begin{classroomqa}
\audiencequestion{ボソンとフェルミオンが同じ質量をもたない場合はどうなるか。\lecturetimestamp{01:01:16}}
\lecturerresponse{その場合、一方を他方に置き換えるとエネルギーが変わるので、提案された超電荷はハミルトニアンと交換せず、単純な超対称性は失われる。
\[
H=m_BN_B+m_FN_F,
\]
に対して、明示的な結果は
\[
[H,Q]=(m_B-m_F)Q.
\]
である。これはちょうど\(m_B=m_F\)のときに消える。\editorialnote{質量が等しくない場合の交換子は、エネルギー変化の議論の直接的な演算子形であり、等質量条件を明示する。}}
\end{classroomqa}

\begin{classroomqa}
\audiencequestion{エネルギー固有状態の議論をもっとゆっくり繰り返せるか。\lecturetimestamp{01:01:37}}
\lecturerresponse{\(H\ket{E}=E\ket{E}\)から始める。質量が等しいことは、\(Q\ket{E}\)が同じ確定したエネルギー\(E\)をもつことを意味するので、\(HQ\ket{E}=EQ\ket{E}\)である。逆の順序では\(QH\ket{E}=QE\ket{E}=EQ\ket{E}\)となる。なぜなら\(E\)は通常のスカラーだからである。それらの差はすべてのエネルギー固有状態で消える。完全性により、この主張はそれらの状態から演算子方程式\([H,Q]=0\)へと高められる。}
\end{classroomqa}

したがって、完全なおもちゃの代数は
\begin{align}
\{Q,\bar Q\}&=H,\\
\{Q,Q\}&=\{\bar Q,\bar Q\}=0,\\
[H,Q]&=[H,\bar Q]=0.
\end{align}
である。超電荷は保存する奇の電荷である。さらに驚くべきことに、それらの代数は粒子の種類を変える操作と時空並進の生成子を混ぜ合わせる。

\section{なぜハミルトニアンが時間並進を意味するのか}

この簡略化された模型には空間座標はないが、時間はある。シュレーディンガー方程式は
\begin{equation}
i\frac{\partial}{\partial t}\ket{\psi(t)}
=H\ket{\psi(t)}.
\end{equation}
と書ける。

\begin{figure}[H]
\centering
\includegraphics[width=0.56\textwidth]{lecture_05_figure_03.png}
\lectureframe{01:07:26.300}
\caption{時間微分は\(\ket{\psi}\)の隣に書かれている。手が分母の一部を隠しているが、完全なシュレーディンガー方程式は上に与えられている。}
\end{figure}

有限のずれ\(\delta\)に対して、
\begin{equation}
\ket{\psi(t+\delta)}
=e^{-iH\delta}\ket{\psi(t)}.
\end{equation}
一次まで展開すると、
\begin{equation}
\ket{\psi(t+\delta)}
=\left(\mathbf 1-i\delta H\right)\ket{\psi(t)}
+O(\delta^2).
\end{equation}
したがって、\(H\)は状態を時間的に更新する無限小生成子である。

アイソスピンやカラーのような内部対称性も粒子のラベルを変える。アイソスピンは陽子と中性子の状態を混ぜることができ、カラー回転は赤、緑、青のクォークを混ぜる。しかし、それらの内部変換のどのような合成も、系を空間的または時間的に並進させない。超対称性は異なる：
\begin{equation}
\begin{gathered}
\text{odd particle-changing transformation}\\
\mathrel{\stackrel{\{\ ,\ \}}{\Longrightarrow}}\\
\text{time translation}
\end{gathered}
\end{equation}
これが、奇の生成子自身が通常とは異なる座標に沿った並進でありうるという手がかりである。

\section{超空間}

\subsection{時間と奇の座標の関数}

関数
\begin{equation}
\Psi=\Psi(\theta,\bar\theta,t).
\end{equation}
を導入する。これはおもちゃの模型における波動関数とみなすこともでき、より一般には超場とみなすこともできる。奇の展開は有限で終わるので、
\begin{equation}
\Psi(\theta,\bar\theta,t)
=A(t)+\bar\chi(t)\theta
+\bar\theta\chi(t)+F(t)\bar\theta\theta.
\end{equation}
この単一の対象は偶と奇の成分関数をまとめる。\editorialnote{この成分展開は先に展開したグラスマン多項式を適用したものであり、おもちゃの超空間上の関数の標準的な明示的内容である。}

通常の時間並進は\(\theta,\bar\theta\)を固定したまま
\begin{equation}
t\longmapsto t+\delta.
\end{equation}
と送る。一次まででは、
\begin{equation}
\Psi(\theta,\bar\theta,t+\delta)
=\Psi(\theta,\bar\theta,t)
+\delta\,\partial_t\Psi(\theta,\bar\theta,t).
\end{equation}

\begin{figure}[H]
\centering
\includegraphics[width=0.72\textwidth]{lecture_05_figure_04.png}
\lectureframe{01:13:15.060}
\caption{黒板では\(t\to t+\delta\)が\(\Psi\)の一次変化の上に置かれ、微分項が直接示されている。}
\end{figure}

\subsection{奇の方向への並進}

次に、奇の座標をずらす変換を考案する：
\begin{equation}
\theta\longmapsto\theta+\epsilon,
\end{equation}
ここで\(\epsilon\)は別の奇のグラスマン・パラメータである。同時に、時間に偶の冪零変位を与える：
\begin{equation}
t\longmapsto t+i\bar\theta\epsilon.
\end{equation}
\(\bar\theta\)と\(\epsilon\)はともに奇であるから、それらの積は偶であり、偶の座標\(t\)に矛盾なく加えることができる。

パラメータを左辺に移すと、
\begin{equation}
i\bar\theta\epsilon=-i\epsilon\bar\theta.
\end{equation}
したがって誘導される変分は
\begin{align}
\delta_\epsilon\Psi
&=\epsilon\,\partial_\theta\Psi
-i\epsilon\bar\theta\,\partial_t\Psi\\
&=\epsilon q_0\Psi,
\end{align}
ここで
\begin{equation}
q_0=\partial_\theta-i\bar\theta\partial_t.
\end{equation}
この奇妙な微分演算子は、合成されたシフトの生成子である。左微分を用いた、整合の取れた共役生成子は
\begin{equation}
\bar q_0=-\partial_{\bar\theta}
+i\theta\partial_t.
\end{equation}

正確な符号は、左グラスマン微分と右グラスマン微分のどちらを用いるか、また共役が奇因子をどのように反転させるかに依存する。規約に依存してはならないのは、時間並進に関する閉包である。上記の定義のもとで、
\begin{align}
\{q_0,\bar q_0\}
&=i\{\partial_\theta,\theta\}\partial_t
+i\{\bar\theta,\partial_{\bar\theta}\}\partial_t\\
&=2i\partial_t.
\end{align}
微分-微分の項は消え、奇座標について二次の項も消える。二つの混合項はそれぞれ \(i\partial_t\) に寄与する。

シュレーディンガー描像の状態に対して \(H=i\partial_t\) であるから、
\begin{equation}
\{q_0,\bar q_0\}=2H.
\end{equation}
ここで
\begin{equation}
q=\frac{q_0}{\sqrt2},
\qquad
\bar q=\frac{\bar q_0}{\sqrt2},
\end{equation}
と規格化すれば、
\begin{equation}
\boxed{\{q,\bar q\}=H.}
\end{equation}
\editorialnote{この講義は、不確かな符号を意図的に検討し、その後で因子 2 に気づく。表示された左微分の規約は、その教室での計算を保ちつつ、一つの自己整合的な規格化を確定する。他の標準的な規約では、符号が座標シフトと生成子の間で移される。}

\begin{classroomqa}
\audiencequestion{超空間の反交換子における最後の因子 2 はどこから来るのか。\lecturetimestamp{01:24:21}}
\lecturerresponse{等しい混合寄与が二つある。一方は \(\{\partial_\theta,\theta\}=1\) から来て、他方は \(\{\partial_{\bar\theta},\bar\theta\}=1\) から来る。それぞれが時間微分に掛かる。したがって、規格化されていない生成子は \(2i\partial_t\) を与える。各超電荷に因子 \(1/\sqrt2\) を入れると、先の規格化が回復する。}
\end{classroomqa}

\begin{classroomqa}
\audiencequestion{この構成は \(\epsilon\) が小さいという仮定を用いたのか。\lecturetimestamp{01:24:51}}
\lecturerresponse{用いたのは、通常の数値的な大きさではなく、冪零性である。\(\epsilon^2=0\) であるから、単一の \(\epsilon\) による展開はすべて一次で止まる。それを無限小または小さいと呼ぶのは有用な類推であるが、打ち切りは厳密である。}
\end{classroomqa}

\editorialnote{グラスマンパラメータには通常の正の大きさはない。それらが無限小パラメータとしての役割を果たすのは、冪零性と展開の厳密な終了による。}

\subsection{幾何学的な意味}

この代数は今や二つの等価な実現を持つ。振動子描像では、\(Q\) と \(\bar Q\) がフェルミオンをボソンに、ボソンをフェルミオンに置き換える。幾何学的描像では、\(q\) と \(\bar q\) が関数を奇方向に平行移動させつつ、通常の時間シフトを誘導する。どちらの描像でも、
\begin{equation}
\begin{gathered}
\{\text{supercharge},\text{conjugate supercharge}\}\\
=\text{time-translation generator}.
\end{gathered}
\end{equation}

\(t,\theta,\bar\theta\) を含む一般化座標空間は超空間である。相対論的場の理論では、それは通常の空間座標も含み、超電荷は \(H=P_0\) だけでなく、完全な時空運動量 \(P_\mu\) について閉じる。この構造を曲げると超重力へと向かう。

ある意味では、グラスマン変数はボソンとフェルミオンを効率的に記帳する道具である。別の意味では、閉包関係はもっと深いことを述べている。粒子の種類と時空の幾何学の区別はもはや絶対的ではない。奇方向に二度動くと、通常の方向への運動が生じる。自然がその構造を実験的に実現しているかどうかは、その構造が整合的で、経済的で、驚くほど剛直であるという数学的事実とは別問題である。
